\chapter{Literature Review}
\label{chap:found}

\section{Introduction}

This chapter examines literature that is available about post-harvest losses and machine learning
in agriculture. It looks at the major concepts, theories, and empirical evidence that inform the
research, especially on aspects that affect post-harvest losses, predictive modeling methods, and
decision support systems. The purpose of the review is to formulate a clear theoretical background
through the synthesis of existing knowledge and to fill gaps that exist in the literature. With the
critical review of past literature, this chapter underscores that there is a need for
straightforward, data-based solutions that suit the agribusinesses of smallholders in developing
nations, and that is why the current research has a factor to play in the field.


\section{Post-Harvest Losses in Developing Countries}

Post-harvest losses may be defined as the increase or decrease in an agricultural product in terms
of amount and quality between the time of harvest and the time of consumption. These losses include
physical loss, nutritional loss, and loss in market value, especially for perishable commodities
like fruits and vegetables. In the developing world, infrastructural, technological, and managerial
deficiencies are particularly acute, and the extent to which the value chain is prone to
post-harvest losses may range from 30--40 per cent of total output~\cite{rajapaksha2021reducing}.

Post-harvest losses may be divided into two subcategories: quantitative losses, which include
weight or volume loss, and qualitative losses, which include loss in texture, taste, nutritional
value, and appearance. Although quantitative losses cause a direct decrease in supply, qualitative
losses typically affect market price or result in rejection, thus exerting a multiplied economic
effect on smallholder farmers~\cite{wakene2024comprehensive}.

There are a number of associated factors that cause post-harvest losses in developing countries,
and one of the most important is storage conditions. Poor storage facilities such as refrigeration
problems or poor ventilation hasten spoilage and the growth of microbes. For example, insufficient
storage of fruits and vegetables in tropical conditions contributes to the rapid physiological
degradation of the product due to high respiration rates and moisture
loss~\cite{rajapaksha2021reducing}. On the same note, most smallholder farmers usually use ancient
forms of storage which cannot sustain the quality of products over time.

Post-harvest turnover is also increased by transportation and logistical inefficiencies. Mechanical
damage and spoilage are caused by poor road systems, extensive transportation times, and the
absence of temperature-regulated transport systems. Rice producers in Liberia face lower
profitability partly because they lose profits in transportation and after
harvesting~\cite{pluato2024impact}. Loading, unloading, and transit operations may cause physical
damage that decreases the shelf life and marketing ability of produce considerably.

Handling practices are also very significant in defining outcomes after harvest. Poor agricultural
practices such as insufficient knowledge of the right harvesting methods, inappropriate packaging,
and rough handling exert more vulnerability to bruising and contamination. Research conducted on
cucumber crop production has indicated that inappropriate handling practices and minimal access to
training by small-scale farmers are primary contributors to post-harvest
losses~\cite{ukwuaba2024mitigating}. The absence of awareness on best practices for post-harvest
management usually adds to these concerns.

The issue is also exacerbated by environmental and climatic conditions. Microbial growth and enzyme
activity are favored by the presence of high temperatures and humidity levels, as is the case in
most developing areas, thereby promoting faster spoilage. Climatic factors have a significant effect
on tomato losses in Ethiopia, especially where storage and cooling facilities are unavailable during
peak harvesting seasons~\cite{wakene2024comprehensive}. Climate variability also interferes with
supply chains, making it difficult to maintain consistent quality after harvest.

The post-harvest losses have significant economic and social impacts. The economic implications are
that losses cause farmers to have less income and profitability, resulting in their inability to
invest in better technologies or practices. In Liberia, smallholder farmers are an example of how a
major source of lost livelihoods is the inefficiencies that occur during the post-harvest
phase~\cite{pluato2024impact}. Socially, post-harvest losses are known to contribute to food
insecurity due to the lack of healthy food in local markets. Furthermore, when there are
inefficiencies in supply chains, food prices rise, burdening the poor population
disproportionately.

Moreover, post-harvest losses compromise overall agricultural development objectives by limiting
food system efficiency. The losses that happen at different stages reduce economic returns and food
availability, as studies on vegetable supply chains such as amaranthus
show~\cite{oluoch2024causes}. It is thus important to address these issues not only to improve
farmer livelihoods but also to increase food security and sustainability in developing countries.


\section{Factors Influencing Post-Harvest Loss}

The intricate interaction of storage, transportation, handling, and environmental factors is the
cause of post-harvest losses in developing countries. These determinants must be understood in order
to identify key predictors of loss risk, which directly answers Research Question~1 of this study.

\subsection{Storage Conditions}

Storage is very important in the maintenance of quality and shelf life of agricultural produce.
Smallholder farmers in most developing nations do not have access to proper materials like cold
rooms or controlled atmosphere facilities to store their goods. This causes fruits and vegetables
to be kept in ambient temperatures, thereby increasing spoilage. According to
\cite{rajapaksha2021reducing}, insufficient storage facilities such as improper ventilation and
temperature control are an important cause of fruit and vegetable spoilage. In the same manner,
\cite{wakene2024comprehensive} mention that tomatoes stored in poor conditions are actively
degraded physiologically, causing quantitative and qualitative losses. Lack of affordable storage
technologies thus continues to be a significant limitation towards minimizing post-harvest losses.

\subsection{Transportation Logistics}

Another significant cause of post-harvest losses is transportation inefficiency. Lack of
appropriate road networks, long distances to travel, and the absence of refrigerated transportation
systems subject produce to mechanical damage and unfavorable environmental conditions.
\cite{pluato2024impact} show that in Liberia, high losses in transportation significantly decrease
the profitability of rice production among smallholder farmers. Any failure to transport produce to
markets promptly raises the chances of spoilage, mainly for perishable goods. Also, poor packaging
during delivery will usually lead to compression and bruising, consequently lowering product quality
and market value.

\subsection{Handling Practices}

The management of practices at different parts of the supply chain, such as harvesting, sorting, and
packaging, has a significant impact on post-harvest results. Misuse of handling techniques may
result in physical injuries to produce, making it more vulnerable to microbial infection and rapid
breakdown. \cite{ukwuaba2024mitigating} identify smallholder cucumber producers, whose major
problems are improper harvesting practises, rough handling, and non-standardised packaging, as
primary contributors to post-harvest losses. Such obstacles are usually associated with inadequate
training of farmers and ignorance of optimal farming practices.

\subsection{Environmental Variables}

Conditions in the environment, especially temperature and humidity, are decisive factors of
post-harvest loss. Hot weather increases the rate of respiration in produce, resulting in quicker
ripening and ultimate spoilage, while elevated moisture content encourages the development of
microorganisms. \-\cite{wakene2024comprehensive} note that losses of tomatoes in Ethiopia are greatly
caused by fluctuations in temperature and humidity. Similarly, \cite{rajapaksha2021reducing}
observe that climate conditions in tropical regions contribute to increasing post-harvest
degradation because there are no mechanisms of temperature control. Seasonal differences also
affect the pattern of losses, where peak harvest seasons often coincide with weather conditions
that favor spoilage.

\subsection{Cohesive Effect of the Supply Chain}

These factors rarely work in isolation; they interact at various stages throughout the supply chain,
increasing aggregate losses. \cite{oluoch2024causes} explain that in the amaranthus supply chain,
losses are experienced during harvesting and distribution to the market due to the triple interaction
of poor storage, unproductive transport, and inefficient handling. This demonstrates the importance
of holistic approaches to post-harvest management.

Overall, post-harvest losses in developing countries are mostly affected by storage conditions,
logistics involved in transportation, post-harvest handling behaviors, and environmental factors.
The identification and quantification of these variables are essential to establishing predictive
models that benefit smallholder agribusinesses through data-driven decisions.


\section{Machine Learning in Agriculture}

Machine learning (ML) as an agricultural application has been on the rise in the past few years due
to the necessity of obtaining better productivity, efficiency, and sustainability in food systems.
ML methods make it possible to analyze large and intricate datasets, getting patterns and making
predictions to help guide decisions. In underdeveloped economies, where the availability of
state-of-the-art technologies is curtailed by resource requirements, easy-to-understand and
interpretable ML models offer a chance at improving agricultural activity, especially in yield
forecasting and crop management.

Among the most outstanding uses of ML in farming, crop yield forecasting can be mentioned. Yield
can be forecasted with sufficient accuracy to assist planning, manage the supply chain, and secure
food. Conventional ways of estimating yields usually take into consideration historic averages or
manual estimations that fail to account for dynamic environmental and operating variables. ML models
on the other hand can combine various variables like weather, soil characteristics, and agricultural
methods to give better predictions. As pointed out by \cite{jabed2024crop}, machine learning
methods, such as regression models and ensemble methods, have been shown to perform highly when
predicting crop yields in various agricultural environments.

Along with predictive yielding, disease detection and crop health monitoring are other crucial areas
of application of ML. The early detection of plant diseases is critical to minimizing losses and
improving crops. The output of machine learning algorithms, especially those that work with image
data, is highly accurate as the algorithms attempt to find patterns in leaf color, texture, and
shape to detect disease symptoms. \cite{meghraoui2024applied} suggest that ML-based disease
detection systems can offer timely responses, eliminating the necessity to conduct manual inspections
or use expensive specialist expertise.

In addition to particular applications, the greater value of ML in agriculture lies in the fact that
it can facilitate processes of data-driven decision-making. Pattern recognition is one of the major
strengths of ML, in which an algorithm can establish trends and relationships among very large and
complex datasets. This feature is particularly applicable to the agricultural setting, where numerous
interacting characteristics can affect outcomes such as yield and post-harvest quality. According to
\cite{jabed2024crop}, this analytical capability has helped increase the accuracy and trustworthiness
of agricultural forecasts.

Another relevant advantage of ML is its use in decision support systems. ML models have the ability
to convert information into practical suggestions to allow farmers and agribusinesses to make good
decisions. Predictive models can, for example, inform decisions regarding harvest timing, storage
approach, and transportation logistics. \cite{meghraoui2024applied} suggest that the application of
ML to the decision-making process in the agricultural domain will make the process more efficient
and less prone to uncertainty, especially in resource-limited contexts. Notably, simple models such
as logistic regression and decision trees can be used to make sure such systems remain interpretable
and accessible to final users.

Although these are the benefits, there are still difficulties in the implementation of ML in the
agricultural sector. Availability and quality of data are usually poor, especially in developing
countries. Moreover, the complexity of models and their usability must be balanced, as very complex
models might not be useful to smallholder farmers. However, the accumulating volume of literature
shows that ML can achieve great changes in the field of agriculture by providing accurate
predictions and enabling informed decision-making processes. In this research, these insights are
applied to forecast post-harvest loss risk using machine learning methods.
